\section{Aula 6}

Hoje vimos um belíssimo teorema em aula sobre decaimento exponencial
e comportamento
de campo médio de $\Theta(p)$ em $\Z^d$.

\begin{theorem}
	Seja $p_c$ a probabilidade crítica de percolação em $\Z^d$. Então, para todo
	$p < p_c$ existem $c = c(p)$ e $C = C(p)$, tal que
	\[
		\Theta_n(p) = \Pr(\vec{0} \leftrightarrow \partial B_n) \leq Ce^{-cn}.
	\]
	Ademais, para $p > p_c$, temos
	\[
		\Theta(p) = \Pr_p(\vec{0} \leftrightarrow \infty) > \frac{p - p_c}{2}.
	\]
\end{theorem}
\begin{remark}
	O decaimento exponencial trata-se justamente da primeira parte do teorema,
	enquanto a segunda parte chama-se comportamento de campo médio. O Augusto
	disse que esse nome vem da similaridade da percolação em $G(n,p)$.
\end{remark}

A prova que veremos foi encontrada pelo Dominil-Copin e pelo Tassion. Ela é bem
espertinha, mas precisa de umas definições e leminhas.

\begin{definition}
	Dado um conjunto finito $S \subset \Z^d$ com $0 \in S$ definimos
	\[
		\phi_p(S) = \sum_{\substack{(x,y) \in E\\x\in S, y\not \in S}} p
		\Pr(0 \leftrightarrow_S x).
	\]
	Onde $\{0 \leftrightarrow_S x\}$ significa que existe um caminho
	aberto que passa
	somente por vértices de $S$.
\end{definition}

De certa forma, $\phi_p(S)$ é a média do número de arestas abertas
conectando $0$ para fora de $S$.
Com essa definição, já conseguimos um leminha bonitinho.
\begin{lemma}
	Se existe $S \subset \Z^d$ finito com $0 \in S$, tal que $\phi_p(S) < 1$,
	então existe $C,c > 0$ tal que, para $n$ suficientemente grande,
	\[
		\Pr(0 \leftrightarrow \partial B_n) \leq Ce^{-cn}.
	\]
\end{lemma}
\begin{proof}
	Suponha que exista tal $S$ e seja $|S| = k$. Em particular,
	se $\{0 \leftrightarrow \partial B_n\}$, como $0 \in S$, deve existir alguma
	aresta aberta $(x,y)$ saindo de $S$ com $y \leftrightarrow \partial B_n$.
	Além do mais, o caminho de $y$ para a borda não precisa voltar para dentro
	de $S$.
	Portanto, pelo union bound e pela observação anterior,
	\[
		\Pr(0 \leftrightarrow \partial B_n) \leq
		\sum_{\substack{(x,y) \in E\\x\in S, y\not \in S}} \Pr(0 \leftrightarrow_S x
		\; \land \; (x,y) \in (\Z^d)_p \; \square \; y \leftrightarrow \partial B_n).
	\]
	Usando a desigualdade de BK \ref{trm:bk} e notando que a
	distância de $y$ até
	$\partial B_n$ é pelo menos $n-k$, temos que o lado direito é menor que
	\[
		\sum_{\substack{(x,y) \in E\\x\in S, y\not \in S}} p\Pr(0 \leftrightarrow_S x)
		\Pr(0 \leftrightarrow \partial B_{n-k}) = \phi_p(S) \Pr(0
		\leftrightarrow \partial B_{n-k}).
	\]
	Por indução, $\Pr(0 \leftrightarrow \partial B_n) \leq (\phi_p(S))^{n/k}$.
\end{proof}

Com esse leminha, definimos uma nova probabilidade crítica que
mostraremos ser igual a $p_c$.
\begin{definition}
	Seja $\tilde{p}_c \in [0,1]$ o supremo dos $p$ tal que existe $S$
	satisfazendo $\phi_p(S) < 1$.
	Crucialmente, definindo
	\[
		A = \bigcup_{0 \in S \subset \Z^d} \{p \in [0,1] : \phi_p(S) < 1\}
	\]
	temos
	\[
		\tilde{p}_c = \sup A
	\]
	e como $A$ é um conjunto aberto, $\tilde{p}_c \not \in A$.

\end{definition}

Segue imediatamente do lema anterior que se $p < \tilde{p}_c$, temos
decaimento exponencial. Vamos mostrar que para $p > \tilde{p}_c$, vale
a segunda parte do teorema. Notamos que isso implica imediatamente
que $\tilde{p}_c = p_c$.

A graça aqui é que supor que $p > \tilde{p}_c$ nos dá muita
informação sobre $(\Z^d)_p$. Sabemos que para
qualquer $S$ descrito como antes, $\phi_p(S) \geq 1$. Seguindo a
intuição anterior, para
todo $S$ a gente em média consegue escapar.

Como era de se esperar,
vamos provar a segunda parte
desse teorema usando a formula de Russo \ref{trm:russo} e uma igualdade esperta.

\begin{lemma}
	Definimos o conjunto aleatório
	\[
		\cS_n = \{v \in B_{n-1} : v \not \leftrightarrow \partial B_n \}
	\]
	de elementos que não conectam ao bordo de $B_n$, temos que
	\begin{equation}
		\frac{d}{dp} \Theta_n (p) = \frac{1}{p(1-p)} \Ex \phi_p(\cS_n)
	\end{equation}
\end{lemma}
\begin{proof}
	Sabemos por \ref{trm:russo} que
	\[
		\frac{d}{dp} \Theta_n(p) = \sum_{e \in E(B_n)} \Pr(e \text{ pivotal
		para } [0 \leftrightarrow \partial B_n])
	\]
	Como $\{e \text{ pivotal para } A\}$ é independente do valor de $e$,
	podemos pagar
	um $(1-p)$ e pedir que $e$ esteja fechada. Portanto,
	\[
		\frac{d}{dp} \Theta_n(p) = \frac{1}{1-p}\sum_{e \in E(B_n)} \Pr(e
		\text{ fechada e } e \text{ pivotal para } [0 \leftrightarrow \partial B_n]).
	\]
	Note que os eventos anteriores implicam justamente que $[0 \not
	\leftrightarrow \partial B_n]$.
	Agora vamos separar o mundo de acordo com $\cS_n = S$.
	Como $[0 \not \leftrightarrow \partial B_n]$, segue que $0 \in
	\cS_n$. Tomando um union bound nos valores de $\cS_n$
	possíveis, temos que o somatório anterior se expressa como
	\[
		\frac{1}{1-p}\sum_{0 \in S \subset B_{n-1}}
		\sum_{e \in E(B_n)}
		\Pr(e \text{ fechada e } e \text{ pivotal para }
		[0 \leftrightarrow \partial B_n] \text{ e } \cS_n = S)
	\]
	Agora $e$ ser pivotal e fechada com $\cS_n = S$ acontece se e
	somente se $e$ conecta
	um elemento de $S$ para outro fora dele e, além do mais,
	que exista um caminho conectando $0$ até $e$ em $S$. Disso podemos
	reescrever a última soma como
	\[
		\frac{1}{1-p}\sum_{0 \in S \subset B_{n-1}}
		\sum_{\substack{(x,y) \in E(\Z^d)\\ x \in S, y \not \in S}}
		\Pr(0 \leftrightarrow_S x \text{ e } \cS_n = S)
	\]
	Como os eventos dentro da probabilidade são independentes, segue que
	\[
		\Theta_n'(p) =
		\frac{1}{1-p}\sum_{0 \in S \subset B_{n-1}} \Pr(\cS_n = S)
		\sum_{\substack{(x,y) \in E(\Z^d)\\ x \in S, y \not \in S}}
		\Pr(0 \leftrightarrow_S x) = \frac{1}{p(1-p)}\Ex \phi_p(\cS_n).
	\]
\end{proof}

Estamos prontos para provar o teorema.

\begin{proof}[Prova do teorema bonitão]
	Seja $\tilde{p}_c$ definido como antes. Já vimos que decaimento
	exponencial vale para $p < \tilde{p}_c$.
	Agora tome $p > \tilde{p}_c$. Se $\Theta_n(p) \geq 1/2$, vale imediatemente
	que $\Theta_n(p) \geq (p - \tilde{p}_c)/2$. Se não, suponha que
	$\Theta_n(p) < 1/2$, nesse caso,
	com probabilidade maior ou igual a $1/2$, $0 \subset \cS_n$. Como $p
	> \tilde{p}_c$,
	para qualquer $S = \cS_n$, temos $\phi_p(S) \geq 1$. Portanto,
	\[
		\Theta_n'(p) \geq \Ex \phi_p(\cS_n) \geq 1 - \Theta_n(p) \geq 1/2.
	\]
	Integrando, temos o resultado para $\Theta_n$. Levando $n \to
	\infty$, recuperamos
	o resultado total. Note que isso implica que $\tilde{p}_c = p_c$.
\end{proof}

\begin{remark}
	Essa definição é legal pois ela já dá (pela observação de $p_c \not \in A$)
	que $\Ex_{p_c}[|C_0|] = \infty$. Basta notar que
	\[
		\Ex_{p_c}[|C_0|] = \sum_{n \geq 1} \sum_{v \in \partial B_n} \Pr(0
		\leftrightarrow_{B_n} v) = \frac{1}{p_c}\sum_{n} \phi_{p_c}(B_n) = \infty.
	\]
\end{remark}

\begin{remark}
	Com um pouco de esforço, vemos também que implica que dado $x,y$ com
	distância $n$ entre eles,
	$\Pr(x \leftrightarrow y) \geq (Cn)^{-d}$. Tomando $x = 0$, temos
	\begin{align*}
		\Pr(0 \leftrightarrow y) \geq \frac{1}{p_c |\partial
		B_n|}\phi_{p_c}(B_{n-1}) \geq \frac{1}{p_c|\partial B_{n}|}.
	\end{align*}
\end{remark}
